\date{}%15 gennaio 2014% versione mandata

\documentclass[11pt]{article}

\usepackage{amsmath,amsfonts,amssymb,amsthm}
\usepackage{latexsym, xspace, enumerate}
\usepackage[mathscr]{eucal}
\usepackage[all]{xypic}
\usepackage{hyperref}
\usepackage{amscd}
\usepackage{vmargin}

\setmarginsrb{28mm}{10mm}{28mm}{25mm}%
          {10mm}{7mm}{10mm}{10mm}


\newtheorem{theorem}{Theorem}[section]
\newtheorem{problem}[theorem]{Problem}
\newtheorem{pro}[theorem]{Proposition}
\newtheorem{proposition}[theorem]{Proposition}
\newtheorem{lemma}[theorem]{Lemma}
\newtheorem{fact}[theorem]{Fact}
\newtheorem{corollary}[theorem]{Corollary}
\newtheorem{conjecture}[theorem]{Conjecture}
\theoremstyle{definition}
\newtheorem{definition}[theorem]{Definition}
\newtheorem{question}[theorem]{Question}
\newtheorem{example}[theorem]{Example}
\newtheorem{remark}[theorem]{Remark}

\newcommand{\R}{\mathbb R}
\newcommand{\N}{\mathbb N}
\newcommand{\C}{\mathbb C}
\newcommand{\G}{G_\delta}
\newcommand{\Q}{\mathbb Q}
\newcommand{\Z}{\mathbb Z}
\newcommand{\Prm}{\mathbb P}
\newcommand{\T}{\mathbb{T}}
\def\ent{\mathrm{ent}}
\def\depth{\mathrm{depth}}
\def\End{\mathrm{End}}
\def\Aut{\mathrm{Aut}}
\def\Im{\mathrm{Im}}
\def\nub{\mathrm{nub}}
\def\supp{\mathrm{supp}}
\def\J{\mathbb{J}}
\newcommand{\NB}{$\clubsuit$}

\renewcommand{\theequation}{\thesection.\arabic{equation}}
\numberwithin{equation}{section}


\title{Topological entropy for automorphisms of\\totally disconnected locally compact groups}

\author{Anna Giordano Bruno\footnote{Research supported by ``Progetti di Eccellenza 2011/12'' of Fondazione CARIPARO and by INdAM}}

\begin{document}

\maketitle

\begin{abstract}
We give a ``limit-free formula'' simplifying the computation of the topological entropy for topological automorphisms of totally disconnected locally compact groups. This result allows us to extend several basic properties of the topological entropy known to hold for compact groups.
\end{abstract}

\bigskip
\noindent \textrm{\small{Key words: topological entropy, totally disconnected locally compact group\\
2010 AMS Subject Classification: 37B40, 22D05, 22D40, 54H11, 54H20, 54C70.}}


\section{Introduction}

In this paper we consider the topological entropy for topological automorphisms of totally disconnected locally compact groups. 

For a locally compact group $G$ let $\mu$ denote a left Haar measure of $G$. Moreover, it is worth recalling immediately that a totally disconnected locally compact group $G$ has as a local base at $1$ the family $\mathcal B(G)$ of all open compact subgroups of $G$, as proved by van Dantzig in \cite{vD}.

\medskip
In \cite{AKM} Adler, Konheim and McAndrew introduced the topological entropy for continuous selfmaps of compact spaces, while later on Bowen in \cite{B} introduced it for uniformly continuous selfmaps of metric spaces, and this definition was extended to uniformly continuous selfmaps of uniform spaces by Hood in \cite{hood}. In particular, as explained in details in \cite{DG-islam}, for continuous endomorphisms of totally disconnected locally compact groups the topological entropy can be introduced as follows. 

Let $G$ be a totally disconnected locally compact group and $\phi:G\to G$ a continuous endomorphism. For every $U\in\mathcal B(G)$ and every positive integer $n$, the \emph{$n$-th $\phi$-cotrajectory} of $U$ is $$C_n(\phi,U)=U\cap \phi^{-1}(U)\cap\ldots\cap\phi^{-n+1}(U),$$ and the \emph{$\phi$-cotrajectory} of $U$ is
$$C(\phi,U)=\bigcap_{n=0}^\infty\phi^{-n}(U)=\bigcap_{n=1}^\infty C_n(\phi,U).$$
Note that $C_n(\phi,U)\in\mathcal B(G)$ for every positive integer $n$, while $C(\phi,U)$ is compact but not open in general, and it is the greatest $\phi$-invariant subgroup of $G$ contained in $U$. Moreover, for every positive integer $n$ the index $[U:C_n(\phi,U)]$ is finite as $U$ is compact and $C_n(\phi,U)$ is open.

The \emph{topological entropy of $\phi$ with respect to $U$} is given by the following limit, which is proved to exist in Lemma \ref{logalpha} below,
$$H_{top}(\phi,U)=\lim_{n\to \infty}\frac{\log[U:C_n(\phi,U)]}{n}.$$
The \emph{topological entropy} of $\phi$ is $$h_{top}(\phi)=\sup\{H_{top}(\phi,U):U\in\mathcal B(G)\}.$$

\medskip
In Section \ref{lf-sec} we prove the main result of this paper, that is the following limit-free formula for the computation of the topological entropy: if $G$ is a totally disconnected locally compact group, $\phi:G\to G$ is a topological automorphism and $U\in\mathcal B(G)$, then
\begin{equation}\label{lff}
H_{top}(\phi,U)=\log[\phi(C(\phi^{-1},U)):C(\phi^{-1},U)].
\end{equation}
%\begin{align*}
%H_{top}(\phi,U)&=\log[\phi^{-1}(C(\phi,U)):C(\phi,U)]+\log\Delta(\phi)\\
%&=\log[\phi(C(\phi^{-1},U)):C(\phi^{-1},U)].
%\end{align*}
This result extends to locally compact groups the limit-free formula proved in \cite{DG-limitfree} in the compact case. That formula was inspired by its counterpart for the algebraic entropy claimed by Yuzvinski in \cite{Y} and verified in \cite{DG-limitfree}. Indeed, in the abelian case the algebraic entropy and the topological entropy are strictly related by the Pontryagin duality (see \cite{DG-BT2,DG-BT,W}).

\medskip
Under some necessary assumptions, the limit-free formula in the compact case was given for continuous endomorphisms, so we leave the following open problem.

\begin{problem}
Is it possible to extend the limit-free formula in \eqref{lff} to continuous endomorphisms of locally compact groups? Under which assumptions?
\end{problem}

An interesting application of the limit-free formula in \eqref{lff} is given in \cite{BDG}, where it allows for a comparison of the topological entropy with the scale function $s(-)$ introduced by Willis in \cite{Willis,Willis2} for topological automorphisms $\phi:G\to G$ of totally disconnected locally compact groups $G$. Indeed, it follows from the so-called tidying procedure given in \cite{Willis2} that
$$s(\phi)=\min\{[\phi(C(\phi^{-1},U)):C(\phi^{-1},U)]:U\in\mathcal B(G)\};$$
so clearly $\log s(\phi)\leq h_{top}(\phi)$ in view of \eqref{lff}.
Furthermore, in \cite{BDG} a characterization is given of the topological automorphisms $\phi:G\to G$ of totally disconnected locally compact groups for which the equality $\log s(\phi)=h_{top}(\phi)$ holds.

\smallskip
In Section \ref{prop-sec}, as consequences of the limit-free formula in \eqref{lff}, we verify the basic properties of the topological entropy, namely, the so-called Invariance under conjugation, Monotonicity, Logarithmic law, Continuity with respect to inverse limits and weak Addition Theorem. Indeed, these properties are known to hold for continuous endomorphisms of compact groups (see \cite{AKM,St}), and now we extend them to the case of topological automorphisms of totally disconnected locally compact groups.

\medskip
In \cite{St} the so-called Addition Theorem was proved for continuous endomorphisms $\phi:G\to G$ of compact groups; indeed, if $N$ is a closed normal $\phi$-invariant subgroup of $G$ and $\overline \phi:G/N\to G/N$ is the continuous endomorphism induced by $\phi$, then $$h_{top}(\phi)=h_{top}(\phi\restriction_N)+h_{top}(\overline\phi).$$
We leave open the following problem to know whether the Addition Theorem holds in the setting of this paper; one should start studying the abelian case.

\begin{problem}
Let $G$ be a totally disconnected locally compact group, $\phi:G\to G$ a topological automorphism, $N$ a closed normal subgroup of $G$ such that $\phi(N)=N$, and $\overline\phi:G/N\to G/N$ the topological automorphism induced by $\phi$. Is it true that $h_{top}(\phi)=h_{top}(\phi\restriction_N)+h_{top}(\overline\phi)$? 
\end{problem}



\section{Limit-free computation of topological entropy}\label{lf-sec}

For a group $G$ and a subgroup $H$ of $G$, we denote by $G/H$ the set of all (left) cosets of $H$ in $G$, and we denote by $[G:H]$ the index of $H$ in $G$, that is the size of $G/H$. If $K$ is another subgroup of $G$, then we denote by $KH/H$ the family of all (left) cosets of $H$ in $G$ with representing elements in $K$, that is $KH/H=\{kH:k\in K\}$, and we denote by $[KH:H]$ the size of this family, generalizing the usual notion of index. If $H\subseteq K$ then we write simply $K/H$ and $[K:H]$, as usual.

We use several times the following easy-to-check properties. %note that only in item (b) we use the generalization of the concept of index.

\begin{fact}\label{gi}
Let $G$ be a group and let $H,K,L$ be subgroups of $G$, with $H\subseteq K$. Then:
\begin{itemize}
\item[(a)] $[G:H]=[G:K][K:H]$;
\item[(b)] $[LH:H]=[L:H\cap L]$;
\item[(c)] $[K:H]\geq[K\cap L:H\cap L]$.
\end{itemize}
If $N$ is a normal subgroup of $G$, $q:G\to G/N$ is the canonical projection and $N\subseteq H$, then
\begin{itemize}
\item[(d)] $[K:H]=[q(K):q(H)]$.
\end{itemize}
\end{fact}

The following is the counterpart for the topological entropy of \cite[Lemma 1.1]{DGSZ} and \cite[Lemma 2.2]{DGS}; its proof follows the one of \cite[Lemma 3.1]{DG-limitfree}. %Its proof follows the one of \cite[Lemma 2.2]{DGS}.

\begin{lemma}\label{logalpha}
Let $G$ be a locally compact group, $\phi:G\to G$ a topological automorphism and $U\in\mathcal B(G)$. For every positive integer $n$ let $c_n:=[U:C_n(\phi,U)]$. Then:
\begin{itemize}
\item[(a)] $c_{n}$ divides $c_{n+1}$ for every $n>0$. 
\end{itemize}
For every $n>0$ let $\alpha_{n}=\frac{c_{n+1}}{c_{n}}=[C_{n}(\phi,U):C_{n+1}(\phi,U)]$. Then:
\begin{itemize}
\item[(b)] $\alpha_{n+1}\leq\alpha_{n}$ for every $n>0$;
\item[(c)] the sequence $\{\alpha_n\}_{n>0}$ stabilizes (i.e., there exist integers $n_0>0$ and $\alpha>0$ such that $\alpha_n=\alpha$ for every $n\geq n_0$);
\item[(d)] $H_{top}(\phi,U)=\log\alpha$.
\end{itemize}
\end{lemma}
\begin{proof}
Let $n>0$. Since there is no possibility of confusion, in this proof we denote $C_n(\phi,U)$ simply by $C_n$.

\smallskip
(a) Since $U\supseteq C_n\supseteq C_{n+1}$ it follows from Fact \ref{gi}(a) that $[U:C_{n+1}]=[U:C_n][C_n:C_{n+1}]$, and so $\frac{c_{n+1}}{c_n}=[C_{n}:C_{n+1}]$; in particular $c_{n}$ divides $c_{n+1}$.

\smallskip
(b) Since $C_{n+1}=C_n\cap \phi^{-1}(U)$, and using Fact \ref{gi}(b), we have $$[C_n:C_{n+1}]=[C_n:C_n\cap \phi^{-n}(U)]=[C_n\cdot\phi^{-n}(U):\phi^{-n}(U)].$$ Now $C_n\subseteq \phi^{-1}(C_{n-1})$, so $$[C_n\cdot\phi^{-n}(U):\phi^{-n}(U)]\leq[\phi^{-1}(C_{n-1})\cdot\phi^{-n}(U):\phi^{-n}(U)].$$
The map $G/\phi^{-n}(U)\to G/\phi^{-n+1}(U)$ induced by $\phi$ is injective, therefore the family of cosets
$(\phi^{-1}(C_{n-1})\cdot\phi^{-n}(U))/\phi^{-n}(U)$ has the same size of its image $(C_{n-1}\cdot\phi^{-n+1}(U))/\phi^{-n+1}(U)$, and so, applying also Fact \ref{gi}(b) and the fact that $C_n=C_{n-1}\cap \phi^{-n+1}(U)$, we have
\begin{equation*}\begin{split}
[\phi^{-1}(C_{n-1})\cdot\phi^{-n}(U):\phi^{-n}(U)]&=[C_{n-1}\cdot\phi^{-n+1}(U):\phi^{-n+1}(U)]=\\&=[C_{n-1}:C_{n-1}\cap \phi^{-n+1}(U)]=[C_{n-1}:{C_n}].
\end{split}\end{equation*}

\smallskip
(c) follows immediately from (b).

\smallskip
(d) By item (c) for $n_0>0$ we have $c_{n_0+n}=\alpha^n c_{n_0}$ for every $n\geq0$, and by the definition of topological entropy 
$$H_{top}(\phi,U)=\lim_{n\to\infty}\frac{\log c_n}{n}=\lim_{n\to \infty}\frac{\log(\alpha^n c_{n_0})}{n}=\log\alpha,$$
and this concludes the proof.
\end{proof}

Lemma \ref{logalpha} yields that the limit defining the topological entropy exists, and gives a precise description for the value of the topological entropy $H_{top}(\phi,U)$ of a topological automorphism $\phi:G\to G$ of a totally disconnected locally compact group $G$ with respect to $U\in\mathcal B(G)$; in particular, this value is in $\log\N_+$, so the value $h_{top}(\phi)\in\log\N_+\cup\{\infty\}$. %In other words, we have the following consequence.
%
%\begin{corollary}
%Consider $h_{top}$ as a map which associates to a pair $(G,\phi)$ of a totally disconnected locally compact group $G$ and a topological automorphism $\phi:G\to G$ the value $h_{top}(\phi)$. Then $h_{top}$ takes values in $\log\N_+\cup\{\infty\}$, so it is a discrete invariant.
%\end{corollary}

\medskip
We are now in position to prove Theorem \ref{lf+}, that is, the main result of this paper, stated in \eqref{lff} in the Introduction. The following folklore fact is needed in the proof. %that we give with a proof for reader's convenience.

\begin{lemma}\emph{\cite{DG-limitfree}}\label{intersect}
Let $G$ be a topological group and let $T$ be a closed subset of $G$. Then for every descending chain 
$B_1 \supseteq B_2 \supseteq \ldots \supseteq B_n \supseteq  \ldots$ of closed subsets of $G$ the intersection $B = \bigcap_{n=1}^\infty B_n $ is non-empty and 
 $\bigcap_{n=1}^\infty (B_{n} T)= B T$, whenever $B_1$ is countably compact. 
\end{lemma}
%\begin{proof} 
%That $B\ne \emptyset$ is a direct consequence of the countable compactness of $B_1$. 
%
%The inclusion $\bigcap_{n=1}^\infty (B_{n} T)\supseteq B T$ is obvious. To verify the converse inclusion pick an element $x\in \bigcap_{n=1}^\infty (B_{n} T)$. Then there exist elements
%$b_n\in B_n$, $t_n \in T$ such that $x = b_nt_n$  for every $n>0$. Let $D_n$ be the closure of the set $\{b_n, b_{n+1}, \ldots \}$ for $n>0$. Then 
%\begin{equation}\label{Eq_19Maggio}
%D_n \subseteq B_n\  \  \mbox{ for each }\  \ n>0. 
%\end{equation}
%The countable compactness of $B_1$ yields that $\bigcap_{n=1}^\infty D_n\ne \emptyset$. Fix an element $b$ of this intersection and note that 
%$b \in B$ due to (\ref{Eq_19Maggio}). It suffices to prove that $b^{-1}x \in T$. Since $T$ is closed it suffices to check that $b^{-1}x$ belongs to the closure of $T$. 
%To this end let $V= V^{-1}$ be a symmetric neighborhood of the neutral element of $G$. Then $bV$ is a neighborhood of $b\in D_1$, so $bV\ni b_m$ for some $m >0$. 
%This yields $Vb^{-1}  \ni b_m^{-1}$, and consequently $Vb^{-1}x  \ni b_m^{-1}x = t_m$. Therefore $Vb^{-1} x\cap T \ne \emptyset$, and so $b^{-1}x$ belongs to the closure of $T$. 
%\end{proof}

%We can now prove the Topological Formula. The following claim is needed in a crucial point of the proof.
%
%\begin{claim}\label{intersect}
%Let $K$ be a compact group and $U$ an open normal subgroup of $G$. Let $\{C_n\}_{n>0}$ be a decreasing chain of compact normal subgroups of $K$.
%Then $\bigcap_{n>0}(C_nU)=(\bigcap_{n>0}C_n)U$. 
%\end{claim}
%\begin{proof}
%The inclusion $\bigcap_{n>0}(C_nU)\supseteq (\bigcap_{n>0}C_n)U$ is obvious. So let $x\in \bigcap_{n>0}(C_nU)$. For every $n>0$ there exist $c_n\in C_n$ and $u_n\in U$ such that $x=c_nu_n$. Since $K$ is compact the sequence $\{c_n\}_{n>0}$ has an accumulation point $c$ and $c\in \bigcap_{n>0}C_n$ because $\{C_n\}_{n>0}$ is a decreasing chain. In particular $c^{-1}x$ is an accumulation point of the sequence $\{c_n^{-1}x\}_{n>0}$ and $c_n^{-1}x=u_n\in U$ for every $n>0$. Hence $c^{-1}x\in U$, and so $x\in cU\subseteq (\bigcap_{n>0}C_n)U$.
%\end{proof}

%The following theorem was proved in \cite[Topological Formula]{DG-limitfree} for continuous endomorphisms of totally disconnected compact groups.

\begin{theorem}\label{lf+}
Let $G$ be a totally disconnected locally compact group, $\phi:G\to G$ a topological automorphism and $U\in\mathcal B(G)$. Then 
$$H_{top}(\phi,U)=\log[\phi(C(\phi^{-1},U)):C(\phi^{-1},U)].$$
\end{theorem}
\begin{proof}
By Lemma \ref{logalpha} there exist integers $n_0>0$ and $\alpha> 0$ such that $\alpha_n=\alpha$ for every $n\geq n_0$, where $\alpha_{n}=[C_{n}(\phi,U):C_{n+1}(\phi,U)]$ for every $n>0$, and $H_{top}(\phi,U)=\log\alpha$. So it suffices to prove that 
\begin{equation}\label{aim}
[\phi(C(\phi^{-1},U)):C(\phi^{-1},U)]=\alpha.
\end{equation}

Since $C(\phi^{-1},U)=U\cap \phi(C(\phi^{-1},U))$, and by Fact \ref{gi}(b), we have
\begin{equation}\label{eqa}\begin{split}
[\phi(C(\phi^{-1},U)):C(\phi^{-1},U)]&=[\phi(C(\phi^{-1},U)):U\cap \phi(C(\phi^{-1},U))]\\&=[\phi(C(\phi,U))\cdot U:U].
\end{split}\end{equation}
The family of cosets $(\phi(U)\cdot U)/U$ is finite, as $U$ is open and $\phi(U)\cdot U$ is compact, since both $U$ and $\phi(U)$ are compact. Consequently, the sequence $\{(\phi(C_n(\phi^{-1},U))\cdot U)/U\}_{n>0}$ is a descending chain of finite subsets of $(\phi(U)\cdot U)/U$ (note that $\phi(C_1(\phi^{-1},U))=\phi(U)$), hence it stabilizes, that is there exists $n_1>0$ such that 
\begin{equation}\label{n1n}
\phi(C_n(\phi^{-1},U))\cdot U=\phi(C_{n_1}(\phi^{-1},U))\cdot U\ \text{for every}\ n\geq n_1;
\end{equation}
in other words, 
\begin{equation}\label{iot}
\bigcap_{n=1}^\infty (\phi(C_{n}(\phi^{-1},U))\cdot U)=\phi(C_{n_1}(\phi^{-1},U))\cdot U.
\end{equation}
Now Lemma \ref{intersect} gives
\begin{equation}\label{lemmagives}
\bigcap_{n=1}^\infty (\phi(C_{n}(\phi^{-1},U))\cdot U)=\left(\bigcap_{n=1}^\infty \phi(C_n(\phi^{-1},U))\right)\cdot U=\phi(C(\phi^{-1},U))\cdot U,
\end{equation}
since $\bigcap_{n=1}^\infty \phi(C_n(\phi^{-1},U))=\phi(C(\phi^{-1},U))$. 
Therefore, \eqref{n1n}, \eqref{iot} and \eqref{lemmagives} yield
\begin{equation}\label{eqb}
\phi(C(\phi^{-1},U))\cdot U=\phi(C_{n_1}(\phi^{-1},U)) \cdot U=\phi(C_{n}(\phi^{-1},U)) \cdot U\ \text{for every}\ n\geq n_1.
\end{equation}
Let $n\geq\max\{n_0,n_1\}$. Then \eqref{eqa} and \eqref{eqb}, together with Fact \ref{gi}(b), give 
\begin{equation}\label{split}\begin{split}
[\phi(C(\phi^{-1},U)):C(\phi^{-1},U)]&=[\phi(C(\phi^{-1},U))\cdot U:U]\\&=[\phi(C_n(\phi^{-1},U))\cdot U:U]\\&=[\phi(C_n(\phi^{-1},U)):\phi(C_n(\phi^{-1},U))\cap U]\\&=[\phi(C_n(\phi^{-1},U)):C_{n+1}(\phi^{-1},U)].
\end{split}\end{equation}
Applying $\phi^{-n}$, we have the equalities 
$$\phi^{-n}(C_{n+1}(\phi^{-1},U))=C_{n+1}(\phi,U)\ \text{and}\ \phi^{-n}(\phi(C_n(\phi^{-1},U)))=C_n(\phi,U),$$ so, since $\phi^{-n}$ is an isomorphism, we have
\begin{equation}\label{lasteq}
[\phi(C_n(\phi^{-1},U)):C_{n+1}(\phi^{-1},U)]=[C_n(\phi,U):C_{n+1}(\phi,U)]=\alpha.
\end{equation}
By \eqref{split} and \eqref{lasteq} we can conclude that \eqref{aim} holds, so the thesis is proved.
%Consequently 
%%\begin{equation}\label{eqc}
%\begin{equation*}\begin{split}
%[\phi(C(\phi^{-1},U)):C(\phi^{-1},U)]=[\phi(C_n(\phi^{-1},U)):C_{n+1}(\phi^{-1},U)]=\\ =\frac{[U:C_{n+1}(\phi^{-1},U)]}{[U:\phi(C_n(\phi^{-1},U))]}=\frac{[U:C_{n+1}(\phi^{-1},U)]}{[\phi(U):C_n(\phi^{-1},U)]}.
%\end{split}\end{equation*}
%\end{equation}
%By Claim \ref{Delta}(b) we have $[\phi(U):C_n]=[U:C_n]\cdot\Delta(\phi)$, therefore recalling that $n\geq n_0$, we can conclude that 
%$$[\phi^{-1}(C(\phi,U)):C(\phi,U)]=\frac{[U:C_{n+1}]}{[U:C_n]}\cdot \Delta(\phi)^{-1}=\alpha\cdot \Delta(\phi)^{-1},$$
%i.e., the wanted equality announced in \eqref{aim}. 
\end{proof}


\section{The topological entropy of the inverse automorphism}

For a locally compact group $G$ and a left Haar measure $\mu$ on $G$, the \emph{modulus} is a group homomorphism $\Delta:\mathrm{Aut(G)}\to\R_+$ such that $\mu(\alpha E)=\Delta(\alpha)\mu(E)$ for every topological automorphism $\alpha\in \Aut(G)$ and every measurable subset $E$ of $G$ (see \cite{HR}). If $G$ is compact, then $\Delta(\phi)=1$ for every topological automorphism $\phi:G\to G$.

\begin{lemma}\label{Delta}
Let $G$ be a totally disconnected locally compact group, $\phi:G\to G$ a topological automorphism and $U\in\mathcal B(G)$. Then
\begin{itemize}
\item[(a)]  $\Delta(\phi)=\frac{\mu(\phi(U))}{\mu(U)}=\frac{[\phi(U):U\cap\phi(U)]}{[U:U\cap\phi(U)]}$.
\item[(b)] If $V\in\mathcal B(G)$ and $V\subseteq U\cap \phi(U)$, then $[\phi(U):V]=[U:V]\cdot\Delta(\phi)$.
\end{itemize}
\end{lemma}
\begin{proof}
(a) Follows from the fact that $\mu(\phi(U))=[\phi(U):U\cap \phi(U)]\cdot\mu(U\cap\phi(U))$ and $\mu(U)=[U:U\cap\phi(U)]\cdot \mu(U\cap\phi(U))$, since $U$ and $\phi(U)$ are compact and $U\cap\phi(U)$ is open, and so $[\phi(U):U\cap\phi(U)]$ and $[U:U\cap\phi(U)]$ are finite.

\smallskip
(b) The situation is described by the following diagram:
$$\xymatrix{
U \ar@{-}[dr]& & \phi(U)\ar@{-}[dl] \\
& U\cap\phi(U)\ar@{-}[d] \\
& V &
}$$
By Fact \ref{gi}(a) we have $$[U:V]=[U:U\cap\phi(U)][U\cap\phi(U):V]$$ and $$[\phi(U):V]=[\phi(U):U\cap\phi(U)][U\cap\phi(U):V].$$ Therefore, $$[\phi(U):V]=[U:V]\frac{[\phi(U):U\cap\phi(U)]}{[U:U\cap\phi(U)]}.$$
To conclude apply item (a).
\end{proof}

In the next proposition we give the local relation between the topological entropy of a topological automorphism of a totally disconnected locally compact group and the topological entropy of its inverse. They coincide precisely when the modulus of $\phi$ is $1$. So for example, as known, in the compact case a topological automorphism has the same topological entropy as its inverse.

\begin{proposition}\label{^-1}
Let $G$ be a totally disconnected locally compact group, $\phi:G\to G$ a topological automorphism and $U\in\mathcal B(G)$. Then
$$H_{top}(\phi^{-1},U)= H_{top}(\phi,U)-\log\Delta(\phi).$$
Hence, $$h_{top}(\phi^{-1})=h_{top}(\phi)-\log\Delta(\phi).$$
\end{proposition}
\begin{proof}
For every positive integer $n$ let $c_n=[U:C_n(\phi,U)]$ and $c_n^*=[U:C_n(\phi^{-1},U)]$.
According to items (a)--(c) of Lemma \ref{logalpha} applied to $\phi$ and $\phi^{-1}$ respectively, $H_{top}(\phi,U)=\log\alpha$ and $H_{top}(\phi^{-1},U)=\log \alpha^*$, where $\alpha$ is the value at which stabilizes the sequence $\alpha_n=\frac{c_{n+1}}{c_{n}}$ and $\alpha^*$ is the value at which stabilizes the sequence $\alpha^*_n=\frac{c_{n+1}^*}{c_{n}^*}$. For every $n>0$, since $\phi^n(C_n(\phi,U))=C_n(\phi^{-1},U)$, using that $\phi^n$ is an automorphism, and by Lemma \ref{Delta}(b), we have that
\begin{equation*}\begin{split}
c_n &=[U:C_n(\phi,U)]=[\phi^n(U):\phi^n(C_n(\phi,U))]=\\&=[\phi^n(U):C_n(\phi^{-1},U))=[U:C_n(\phi^{-1},U)]\cdot\Delta(\phi^n)=c_n^*\cdot \Delta(\phi^n).
\end{split}\end{equation*}
Therefore, since $\Delta$ is a homomorphism, $$\alpha=\frac{c_{n+1}}{c_n}=\frac{c_{n+1}^*}{c_n^*}\cdot\frac{\Delta(\phi^{n+1})}{\Delta(\phi^n)}=\alpha^*\Delta(\phi),$$ hence we have the equality $\log\alpha=\log\alpha^*+\log\Delta(\phi)$, which gives the first assertion of the proposition.

The second equality in the thesis follows from the first one taking the supremum for $U$ ranging in $\mathcal B(G)$.
\end{proof}

The following is an immediate consequence of Theorem \ref{lf+} and Proposition \ref{^-1}.

\begin{corollary}\label{lf-}
Let $G$ be a totally disconnected locally compact group, $\phi:G\to G$ a topological automorphism and $U\in\mathcal B(G)$. Then 
$$H_{top}(\phi,U)=\log[\phi^{-1}(C(\phi,U)):C(\phi,U)]+\log\Delta(\phi).$$
\end{corollary}
\begin{proof}
Theorem \ref{lf+} applied to $\phi^{-1}$ gives $$H_{top}(\phi^{-1},U)=\log[\phi(C(\phi^{-1},U)):C(\phi^{-1},U)].$$ Moreover, Proposition \ref{^-1} yields $H_{top}(\phi,U)= H_{top}(\phi^{-1},U)+\log\Delta(\phi)$, hence the thesis follows.
\end{proof}

%We give a direct proof of Corollary \ref{lf-}.
%
%\begin{proof}[\bf Second proof of Corollary \ref{lf-}.]
%By Lemma \ref{logalpha} there exist integers $n_0>0$ and $\alpha> 0$ such that $\alpha_n=\alpha$ for every $n\geq n_0$ and $H_{top}(\phi,U)=\log\alpha$. So it suffices to prove that 
%\begin{equation}\label{aim}
%[\phi^{-1}(C(\phi,U)):C(\phi,U)]=\alpha\cdot \Delta(\phi)^{-1}.
%\end{equation}
%We start noting that 
%\begin{equation}\label{eqa}
%[\phi^{-1}(C(\phi,U)):C(\phi,U)]=[\phi^{-1}(C(\phi,U)):\phi^{-1}(C(\phi,U))\cap U]=[\phi^{-1}(C(\phi,U))\cdot U:U].
%\end{equation}
%The family of cosets $(\phi^{-1}(U)\cdot U)/U$ is finite, as $U$ is open and $\phi^{-1}(U)$ is compact, and $\{(\phi^{-1}(C_n)\cdot U)/U\}_{n>0}$ is a descending chain of finite subsets, hence it stabilizes, that is there exists $n_1>0$ such that $$\phi^{-1}(C_n)\cdot U=\phi^{-1}(C_{n_1})\cdot U\ \text{for every}\ n\geq n_1;$$ in other words $\bigcap_{n=1}^\infty (\phi^{-1}(C_{n})\cdot U)=\phi^{-1}(C_{n_1})\cdot U$. Lemma \ref{intersect} gives
% $$
% \bigcap_{n=1}^\infty (\phi^{-1}(C_{n})\cdot U)=\left(\bigcap_{n=1}^\infty \phi^{-1}(C_n)\right)\cdot U,
% $$
%and $\bigcap_{n=1}^\infty \phi^{-1}(C_n)=\phi^{-1}(C(\phi,U))$,  therefore 
%\begin{equation}\label{eqb}
%\phi^{-1}(C(\phi,U))\cdot U=\phi^{-1}(C_n) \cdot U\ \text{for every}\ n\geq n_1.
%\end{equation}
%Let $n\geq\max\{n_0,n_1\}$. Then \eqref{eqa} and \eqref{eqb} give 
%\begin{equation*}\begin{split}
%[\phi^{-1}(C(\phi,U)):C(\phi,U)]=[\phi^{-1}(C(\phi,U))\cdot U:U]=\\=[\phi^{-1}(C_n)\cdot U:U]=[\phi^{-1}(C_n):\phi^{-1}(C_n)\cap U]=[\phi^{-1}(C_n):C_{n+1}].
%\end{split}\end{equation*}
%Consequently 
%%\begin{equation}\label{eqc}
%$$[\phi^{-1}(C(\phi,U)):C(\phi,U)]=[\phi^{-1}(C_n):C_{n+1}]=\frac{[U:C_{n+1}]}{[U:\phi^{-1}(C_n)]}=\frac{[U:C_{n+1}]}{[\phi(U):C_n]}.$$
%%\end{equation}
%By  \ref{Delta}(b) we have $[\phi(U):C_n]=[U:C_n]\cdot\Delta(\phi)$, therefore recalling that $n\geq n_0$, we can conclude that 
%$$[\phi^{-1}(C(\phi,U)):C(\phi,U)]=\frac{[U:C_{n+1}]}{[U:C_n]}\cdot \Delta(\phi)^{-1}=\alpha\cdot \Delta(\phi)^{-1},$$
%i.e., the wanted equality announced in \eqref{aim}. 
%\end{proof}


\section{Properties of topological entropy}\label{prop-sec}

In this section we consider the basic properties of the topological entropy. When the group is compact it is known that they hold, but thanks to the limit-free formula proved in Theorem \ref{lf+} it is possible to extend them to the case of topological automorphisms of totally disconnected locally compact groups.

\medskip
We start proving in our setting that the function $H_{top}(\phi,-)$ is antimonotone. It is already known from \cite{DG-islam}.

\begin{lemma}\label{antim}
Let $G$ be a totally disconnected locally compact groupa and $\phi:G\to G$ a continuous endomorphism.
For any $U,V\in\mathcal B(K)$, if $U\subseteq V$, then $H_{top}(\phi,V)\leq H_{top}(\phi,U)$.
\end{lemma}
\begin{proof}
Since $U\subseteq V$, we have that $C_n(\phi,U)\subseteq C_n(\phi,V)$ for every $n>0$. Then, also by Fact \ref{gi}(a) we have for every $n>0$
$$[V:C_n(\phi,V)]\leq[V:C_n(\phi,U)]=[V:U][U:C_n(\phi,U)].$$
Applying the definition of topological entropy, that is, taking the logarithm, dividing by $n$ and letting $n\to\infty$, we find that $H_{top}(\phi,V)\leq H_{top}(\phi,U)$.
\end{proof}

By the definition of topological entropy we immediately derive from Lemma \ref{antim} that, in order to compute the ``global'' topological entropy, it suffices to take the supremum when $U$ ranges in a local base at $1$ of $K$:

\begin{corollary}\label{basesuff}
Let $G$ be a totally disconnected locally compact group, $\phi:G\to G$ a continuous endomorphism and $\mathcal B\subseteq \mathcal B(G)$ a local base at $1$. Then $h_{top}(\phi)=\sup\{H_{top}(\phi,U):U\in\mathcal B\}$.
\end{corollary}

We start proving that the topological entropy is invariant under conjugation.

\begin{proposition}\label{conj}
Let $G$ be a totally disconnected locally compact group and $\phi:G\to G$ a topological automorphism. Let $H$ be another totally disconnected locally compact group and $\xi:G\to H$ a topological isomorphism. Then $\xi(C(\phi,U))=C(\xi\phi\xi^{-1},\xi(U))$ for every $U\in\mathcal B(G)$, and so $h_{top}(\phi)=h_{top}(\xi\phi\xi^{-1})$.
\end{proposition}
\begin{proof}
Let $\psi=\xi\phi\xi^{-1}$. It is clear that, for $U\in\mathcal B(G)$,
\begin{equation}\label{uff}
\xi(C(\phi,U))=C(\psi,\xi(U)).
\end{equation}

Since $\xi:G\to H$ is a topological isomorphism, it induces a bijection $\mathcal B(G)\to \mathcal B(H)$, so by Theorem \ref{lf+} it suffices to check that 
\begin{equation}\label{suff}
[\phi(C(\phi^{-1},U)):C(\phi^{-1},U)]=[\psi(C(\psi^{-1},\xi(U))):C(\psi^{-1},\xi(U))].
\end{equation}
From \eqref{uff} applied to $\phi^{-1}$ and $\psi^{-1}$ we have $C(\psi^{-1},\xi(U))=\xi(C(\phi^{-1},U))$, so, using also the fact that $\xi$ is an isomorphism,
\begin{equation*}\begin{split}
[\psi(C(\psi^{-1},\xi(U))):C(\psi^{-1},\xi(U))]=[\psi(\xi(C(\phi^{-1},U))):\xi(C(\phi^{-1},U))]=\\
=[\xi(\phi(C(\phi^{-1},U))):\xi(C(\phi^{-1},U))]=[\phi(C(\phi^{-1},U)):C(\phi^{-1},U)];
\end{split}\end{equation*}
then \eqref{suff} is verified and this concludes the proof.
\end{proof}

In the next proposition we prove that the monotonicity of the topological entropy under taking subgroups and quotients.

%\begin{lemma}\label{xiU+}
%Let $G$ be a totally disconnected locally compact group, $\phi:G\to G$ a topological automorphism and $U\in\mathcal B(G)$. 
%\begin{itemize}
%\item[(a)] If $N$ is a normal closed subgroup of $G$ such that $\phi(N)=N$, then $C(\phi\restriction_N,U\cap N)=C(\phi,U)\cap N$;
%\item[(b)] if $q:G\to G/N$ is the canonical projection, then $C(\overline\phi,q(U))=q(C(\phi,U))$;
%\end{itemize}
%\end{lemma}

\begin{proposition}\label{monotonicity}
Let $G$ be a totally disconnected locally compact group, $\phi:G\to G$ a topological automorphism and $N$ a closed normal subgroup of $G$ such that $\phi(N)=N$, and let $\overline\phi:G/N\to G/N$ be the topological automorphism induced by $\phi$. Then:
\begin{itemize}
\item[(a)] $C(\phi\restriction_N,U\cap N)=C(\phi,U)\cap N$ for every $U\in\mathcal B(G)$, and $h_{top}(\phi)\geq h_{top}(\phi\restriction_N)$;
\item[(b)] $C(\overline\phi,q(U))=q(C(\phi,U))$ for every $U\in\mathcal B(G)$ such that $N\subseteq U$, where $q:G\to G/N$ is the canonical projection; moreover, $h_{top}(\phi)\geq h_{top}(\overline\phi)$.
\end{itemize}
\end{proposition}
\begin{proof}
(a) The first assertion is clear. Let $V\in \mathcal B(N)$. Then $V= U\cap N$ for some $U\in\mathcal B(G)$. We verify that $$H_{top}(\phi\restriction_N,V)\leq H_{top}(\phi,U)$$ and this implies the wanted inequality.
By the first assertion in item (a), $C((\phi\restriction_N)^{-1},V)=C(\phi^{-1},U)\cap N$, so Fact \ref{gi}(c) gives
\begin{equation*}\begin{split}
[\phi\restriction_N(C((\phi\restriction_N)^{-1},V)):C((\phi\restriction_N)^{-1},V)]&=[\phi(C(\phi^{-1},U)\cap N):C(\phi^{-1},U)\cap N]\\&\leq[\phi(C(\phi^{-1},U)):C(\phi^{-1},U)].
\end{split}\end{equation*}
By Theorem \ref{lf+}, taking the logarithms, this inequality gives $H_{top}(\phi\restriction_N,V)\leq H_{top}(\phi,U)$.

\smallskip
(b) The proof of the first assertion is straightforward. Let $q:G\to G/N$ be the canonical projection and $q(U)\in\mathcal B(G/N)$, where $U\in\mathcal B(G)$; we can assume without loss of generality that $N\subseteq U$. Since $N\subseteq C(\phi^{-1},U)$, by the first assertion in item (b) and Fact \ref{gi}(d) we have 
\begin{equation*}\begin{split}
[\overline\phi(C(\overline\phi^{-1},q(U))):C(\overline\phi^{-1},q(U))]=[\overline\phi(q(C(\phi^{-1},U))):q(C(\phi^{-1},U))]=\\
=[q(\phi(C(\phi^{-1},U))):q(C(\phi^{-1},U))]]=[\phi(C(\phi^{-1},U)):C(\phi^{-1},U)].
\end{split}\end{equation*}
Theorem \ref{lf+} applies to this equality and gives $H_{top}(\overline\phi,q(U))=H_{top}(\phi,U)$, hence we can conclude that $h_{top}(\overline\phi)\leq h_{top}(\phi)$.
\end{proof}

We go on proving the so-called Logarithmic law for the topological entropy. Note that in the compact case $h_{top}(\phi)=h_{top}(\phi^{-1})$ and so the Logarithmic law extends to all integers. This is not possible in general in the locally compact case, as follows immediately from Proposition \ref{^-1}, where the modulus of $\phi$ is involved.

\begin{proposition}
Let $G$ be a totally disconnected locally compact group, $\phi:G\to G$ a topological automorphism and $k$ a positive integer. Then $h_{top}(\phi^k)=k\cdot h_{top}(\phi)$.
\end{proposition}
\begin{proof}
Let $U\in\mathcal B(G)$ and denote $C=C(\phi^{-1},U)$. We start noting that $$C=C(\phi^{-k},V),\ \text{where}\ V=C_k(\phi^{-1},U).$$
 Moreover, $\phi^{n+1}(C)\supseteq \phi^n(C)$ and $[\phi^{n+1}(C):\phi^n(C)]=[\phi(C):C]$ for every $n\geq0$, applying the automorphism $\phi^{-n}$. Then Fact \ref{gi}(a) gives 
\begin{align*}
[\phi^k(C):C]&=[\phi^k(C):\phi^{k-1}(C)]\cdot[\phi^{k-1}(C):\phi^{k-2}(C)]\cdot\ldots\cdot[\phi(C):C]\\
&=[\phi(C):C]^k.
\end{align*}
Therefore, $$[\phi^k(C(\phi^{-k},V)):C(\phi^{-k},V)]=[\phi^k(C):C]=[\phi(C):C]^k;$$
so Theorem \ref{lf+} implies $$H_{top}(\phi^k,V)=k\cdot H_{top}(\phi,U).$$
This immediately gives $h_{top}(\phi^k)\geq k\cdot h_{top}(\phi)$.
By Lemma \ref{antim}, since $V\subseteq U$, we have also $$H_{top}(\phi^k,U)\leq H_{top}(\phi^k,V)=k \cdot H_{top}(\phi,U),$$
so the missing converse inequality $h_{top}(\phi^k)\leq k\cdot h_{top}(\phi)$ holds as well, and this concludes the proof.
\end{proof}

We verify now the continuity of the topological entropy with respect to inverse limits.

\begin{proposition}
Let $G$ be a totally disconnected locally compact group and $\phi:G\to G$ a topological automorphism. If $\{N_i:i\in I\}$ is a directed system of closed normal subgroups of $G$ with $\phi(N_i)=N_i$ and $\bigcap_{i\in I} N_i=\{1\}$, then $G\cong \varprojlim G/N_i$ and $h_{top}(\phi)=\sup_{i\in I} h_{top}(\overline \phi_i)$, where $\overline \phi_i:G/N_i\to G/N_i$ is the continuous endomorphism induced by $\phi$.
\end{proposition}
\begin{proof}
By Proposition \ref{monotonicity}(b) the inequality $h_{top}(\phi)\geq\sup_{i\in I}h_{top}(\overline\phi_i)$ holds. So let $U\in\mathcal B(G)$. There exists $i\in I$ such that $N_i\subseteq U$. Let $q_i:G\to G/N_i$ be the canonical projection. Since $N_i\subseteq C(\phi^{-1},U)$, by Proposition \ref{monotonicity}(b) and Fact \ref{gi}(d) we have 
\begin{equation*}\begin{split}
[\overline\phi_i(C(\overline\phi_i^{-1},q_i(U))):C(\overline\phi_i^{-1},q_i(U))]=[\overline\phi_i(q_i(C(\phi^{-1},U))):q_i(C(\phi^{-1},U))]\\
=[q_i(\phi(C(\phi^{-1},U))):q_i(C(\phi^{-1},U))]=[\phi(C(\phi^{-1},U)):C(\phi^{-1},U)].
\end{split}\end{equation*}
Theorem \ref{lf+} applies to this equality and gives $H_{top}(\overline\phi_i,q_i(U))=H_{top}(\phi,U)$. Therefore $H_{top}(\phi,U)\leq h_{top}(\overline\phi_i)$, hence we can conclude that $h_{top}(\phi)\leq\sup_{i\in I}h_{top}(\overline\phi_i)$.
\end{proof}

The following property is the weak Addition Theorem for the topological entropy.

\begin{theorem}
Let $G$ and $H$ be totally disconnected locally compact groups, $\phi:G\to G$ and $\psi:H\to H$ topological automorphisms. Then $h_{top}(\phi\times \psi)=h_{top}(\phi)+h_{top}(\psi)$.
\end{theorem}
\begin{proof}
Let $\eta=\phi\times\psi$. Since the family $\mathcal B=\mathcal B(G)\times\mathcal B(H)$ is a local base at $1$ of $G\times H$, in view of Corollary \ref{basesuff} it suffices to take $U\in\mathcal B(G)$ and $V\in\mathcal B(H)$, and to verify that $$H_{top}(\eta,U\times V)=H_{top}(\phi,U)+H_{top}(\psi,V);$$
this holds since $C(\eta^{-1},U\times V)=C(\phi^{-1},U)\times C(\psi^{-1},V)$, so
$$[\eta(C(\eta^{-1},U\times V)):C(\eta^{-1},U\times V)]=[\phi(C(\phi^{-1},U)):C(\phi^{-1},U)]\cdot[\psi(C(\psi^{-1},V)):C(\psi^{-1},V)],$$
and Theorem \ref{lf+} concludes the proof.
\end{proof}


%\section{Topological entropy and the scale function}
%
%For $G$ a totally disconnected locally compact group and $\phi:G\to G$ a topological automorphism, it was noted in \cite{CH} that, for any $U\in\mathcal B(G)$,
%$$\log s(\phi)=\lim_{n\to\infty}\frac{\log[\phi^n(U):U\cap \phi^n(U)]}{n}.$$
%This gives immediately that
%\begin{equation}\label{geq}
%h_{top}(\phi)\geq \log s(\phi).
%\end{equation}
%
%We give now an example witnessing that this inequality can be strict.
%
%\begin{example}
%Let $p$ be a prime and $G=\Z(p^\infty)^{\Z}$. Imposing that $U=\Z(p)^\Z$ is open and compact in $G$, then $G$ is given a locally compact (non-compact) topology. Consider $\sigma:G\to G$ the left shift. Clearly $\sigma(U)=U$. Then:
%\begin{itemize}
%\item[(a)] $s(\sigma)=1$;
%\item[(b)] $H_{top}(\sigma,U)=0$;
%\item[(c)] $H_{top}(\sigma,V)=\log p$, where $V=\Z(p)^{-\N_+}\oplus\{0\}\oplus\Z(p)^{\N_+}$; in fact $[\sigma(V_+):V_+]=p$.
%\end{itemize}
%Note that $V_+=\Z(p)^{\N_+}$ and $V_-=\Z(p)^{-\N_+}$, therefore $V$ is tidy above. On the other hand $V_{++}=\Z(p)^{(-\N_+)}\oplus\{0\}\oplus\Z(p)^{\N_+}$, which is dense in $U$ and so it is not closed, in other words $V$ is not tidy below.
%\end{example}
%
%From the tidying procedure it follows that
%$$\log s(\phi)=\min\{\log [\phi(C(\phi^{-1},U)):C(\phi^{-1},U)]:U\in\mathcal B(G)\}.$$
%Furthermore Theorem \ref{lf} gives
%$$h_{top}(\phi)=\sup\{\log[\phi(C(\phi^{-1},U)):C(\phi^{-1},U)]:U\in\mathcal B(G)\}.$$
%Therefore we find again the inequality \eqref{geq}. On the other hand the latter two formulas give the following result.
%
%\begin{proposition}
%Let $G$ be a totally disconnected locally compact group and $\phi:G\to G$ a topological automorphism. If $\nub(\phi)=\{1\}$, then $h_{top}(\phi)=\log s(\phi)$.
%\end{proposition}
%\begin{proof}
%For every tidy $U\in\mathcal B(G)$ we have $\log s(\phi)=\log[\phi(C(\phi^{-1},U)):C(\phi^{-1},U)]$, so $H_{top}(\phi,U)=\log s(\phi)$ by Theorem \ref{lf+}. Since $\nub(\phi)=\{1\}$ precisely when the tidy subgroups form a local base at $1$,  \ref{basesuff} gives $h_{top}(\phi)=\log s(\phi)$.
%\end{proof}




%\section{Bridge Theorem}
%
%
%%In this section we consider topological end algebraic entropy for continuous endomorphisms of locally compact abelian groups, proving a Bridge Theorem for totally disconnected locally compact abelian groups, which extends Weiss Bridge Theorem \ref{Weiss}.
%%
%%\medskip
%%We start from topological entropy.
%%In \cite{B} Bowen defined  the topological entropy for uniformly continuous self-maps of metric spaces, 
%%%and gave an alternative description of its topological entropy 
%%%for uniformly continuous self-maps $T$ of locally compact metric spaces $X$ endowed with a $T$-homogeneous measure. In 1974 Hood \cite{hood} noticed that Bowen's definition of topological entropy, as well as its equivalent description, could be extended respectively to uniformly continuous self-maps of uniform spaces and to uniformly continuous self-maps $T$ of locally compact uniform spaces $X$ with a $T$-homogeneous measure. 
%%and this notion was extended by Hood in \cite{hood} to uniformly continuous self-maps of uniform spaces.
%%%
%%%It is worth mentioning that, for a 
%%%Given continuous self-map $T: X \to X$ of a compact space $X$, consider $X$ with its unique admissible uniform structure $\mathcal U$; then $T:(X,\mathcal U)\to(T,\mathcal U)$ is uniformly continuous and for $T$ the topological entropy 
%%%defined by Hood
%%For continuous self-maps of compact spaces this topological entropy coincides with the one defined by Adler, Konheim and McAndrew (see \cite[Corollary 2.14]{DSV}).
%%%
%%%gave two definitions of topological entropy, the former for a uniformly continuous self-map of a metric space, and the latter for a uniformly continuous self-map $T$ of a locally compact metric space $X$ endowed with a $T$-homogeneous measure. In 1974 Hood \cite{hood} noticed that Bowen's definitions could be extended respectively for a uniformly continuous self-map of a uniform space and for a uniformly continuous self-map $T$ of a locally compact uniform space $X$ endowed with a $T$-homogeneous measure. It is worth mentioning that all the above notions coincide whenever they are defined (see \cite{B}, \cite{DSV} and \cite{hood}).%\{Mettiamo qui anche il libro di Ward?}. 
%%
%%%Consider now an locally compact abelian group $G$ with a Haar measure $\mu$, and an endomorphism $\phi:G\to G$. In particular, $G$ is a locally compact uniform space when endowed with its canonical left uniformity $\mathcal U$; furthermore, $\phi:(G,\mathcal U)\to (G,\mathcal U)$ is uniformly continuous, and $\mu$ is $\phi$-homogeneous. %in the sense of Bowen and Hood. 
%%%Hence, 
%%
%%Moreover, Hood's extension of Bowen's definition of topological entropy applies to continuous endomorphisms $\phi$ of locally compact abelian groups $G$, and can be given in the following way. Denote by $\mathcal C(G)$ the family of compact neighborhoods of $0$ in $G$. For every $U\in\mathcal C(G)$ and every positive integer $n$,
%%$$C_n(\phi,U)=U\cap\phi^{-1}(U)\cap\ldots\cap \phi^{-n+1}(U)$$
%%is the {\em $n$-th $\phi$-cotrajectory of $U$}. The \emph{topological entropy} of $\phi$ with respect to $U$ is
%%$$H_{top}(\phi,U)=\limsup_{n\to\infty}\frac{-\log\mu(C_n(\phi,U))}{n}$$
%%and the \emph{topological entropy} of $\phi$ is
%%$$h_{top}(\phi)=\sup\{H_{top}(\phi,U\}:U\in\mathcal C(G)\}.$$
%%
%%\medskip
%%In \cite{Pet1} Peters gave a further generalization of the entropy he defined in \cite{Pet}, using Haar measure, for topological automorphisms of locally compact abelian groups. 
%%%Let $G$ be an locally compact abelian group with a Haar measure $\mu$, and let $\phi:G\to G$ be an automorphism. Then
%%% {\em Peters entropy} of $\phi$ is
%%%$$h_\infty(\phi)=\left\{\lim_{n\to\infty}\frac{\log\mu(T_n(\phi^{-1},C))}{n}:C\in\c(G)\right\} .$$
%%In \cite{V} Virili modified Peter's definition, in the same way as in \cite{DG2} for the discrete case, obtaining a new notion of algebraic entropy for endomorphisms $\phi$ of locally compact abelian groups $G$.
%%%We indicate it  by $h_{alg}(-)$ as it coincides on endomorphisms of discrete Abelian groups with the already defined algebraic entropy from \cite{DG} (see Example \ref{ent=h_{alg}}(b)). For the precise relation between $h_\infty(-)$ and $h_{alg}(-)$ see Remark \ref{h_infty=h_{alg}}.
%%Let $U\in\mathcal C(G)$. The \emph{algebraic entropy of $\phi$ with respect to $U$} is
%%$$H_{alg}(\phi,U)=\limsup_{n\to\infty}\frac{\log\mu(T_n(\phi,U))}{n},$$
%%and the \emph{algebraic entropy} of $\phi$ is
%%$$h_{alg}(\phi)=\sup\{H_{alg}(\phi,U):U\in\mathcal C(G)\} .$$
%%
%%\bigskip
%%
%%In order to compute the topological and the algebraic entropy, it is a direct consequence of the monotonicity of the functions $H_{top}(\phi,-)$ and $H_{alg}(\phi,-)$ that it suffices to consider respectively a base of the neighborhoods of $0$ in $G$ and a cofinal subset of $\mathcal C(G)$:
%%
%\begin{lemma}\label{base}
%Let $G$ be a locally compact abelian group.
%\begin{itemize}
%\item[(a)] If $\mathcal B\subseteq \mathcal C(G)$ is a base of the neighborhoods of $0$ in $G$, then $h_{top}(\phi)=\sup\{H_{top}(\phi,U\}:U\in\mathcal B\}$.
%\item[(b)] If $\mathcal B\subseteq \mathcal C(G)$ is cofinal, then $h_{alg}(\phi)=\sup\{H_{alg}(\phi,U):U\in\mathcal B\}$.
%\end{itemize}
%\end{lemma}
%%
%%As it is proved in \cite{DG-islam}, for totally disconnected locally compact abelian groups, it is possible to avoid the use of Haar measure, and to prove that $\limsup$ in the definitions are actually limits. More precisely, the limit exists when $U$ is a compact open subgroup. Let $\mathcal B(G)$ be the family of all compact open subgroups of a locally compact abelian group $G$. In particular, $\mathcal B(G)\subseteq \mathcal C(G)$.
%%
%%\begin{proposition}\label{no-mu}
%%Let $G$ be a locally compact abelian group, let $\phi:G\to G$ be a continuous endomorphism and $U\in\mathcal B(G)$. Then:
%%\begin{itemize}
%%\item[(a)] $H_{top}(\phi,U)=\lim_{n\to \infty}\frac{\log[U:C_n(\phi,U)]}{n}$;
%%\item[(b)] $H_{alg}(\phi,U)=\lim_{n\to\infty}\frac{\log[T_n(\phi,U):U]}{n}$.
%%\end{itemize}
%%\end{proposition}
%%%\begin{proof}
%%%(a)
%%%
%%%\smallskip
%%%(b)
%%%\end{proof}
%
%\begin{lemma}\label{B(G)}
%Let $G$ be a locally compact abelian group. 
%\begin{itemize}
%\item[(a)]If $U\in\mathcal B(G)$, then $U^\perp\in\mathcal B(\widehat G)$.
%\item[(b)]If $K\in\mathcal B(\widehat G)$, then $K^\top\in\mathcal B(G)$.
%\end{itemize}
%\end{lemma}
%\begin{proof}
%(a) We have the following corresponding exact sequences:
%$$0\to U\to G \to G/U \to 0$$
%and
%$$0\leftarrow \widehat U\cong G/U^\perp	\leftarrow\widehat G\leftarrow \widehat{G/U}\cong U^\perp\leftarrow 0.$$
%Since $U$ is open in $G$, the quotient $G/U$ is discrete, therefore $U^\perp\cong \widehat{G/U}$ is compact. Moreover, $U$ is compact, so $G/U^\perp\cong\widehat U$ is discrete, and hence $U$ is open in $G$.
%
%\smallskip
%(b) follows from (a).
%\end{proof}
%
%\begin{lemma}\label{(v)}
%Let $G$ be a locally compact abelian group.
%If $U_1,\ldots,U_n$ are compact open subgroups of $G$, then $(\sum_{i=1}^nU_i)^\perp\cong \bigcap_{i=1}^nU_i^\perp$ and $(\bigcap_{i=1}^nU_i)^\perp\cong\sum_{i=1}^n U_i^\perp$.
%\end{lemma}
%
%\begin{lemma}\label{perp}
%Let $G$ be a locally compact abelian group, let $\phi:G\to G$ be a continuous endomorphism and $U\in\mathcal B(G)$.
%For every $n\in\N$, $$(\phi^{-n}(U))^\perp=(\widehat\phi)^n (U^\perp).$$
%Then, for every $n\in\N_+$,  $$C_n(\phi,U)^\perp=T_n(\widehat\phi,U^\perp).$$
%\end{lemma}
%\begin{proof}
% We prove the result for $n=1$, that is, $(\phi^{-1}(U))^\perp=\widehat\phi (U^\perp)$. The proof for $n>1$ follows easily from this case noting that $(\widehat\phi)^n=\widehat{(\phi^n)}$. 
%
%Let $\pi':G\to G/\phi^{-1}(U)$ and $\pi:G\to G/U$ be the canonical projections. Let $\widetilde\phi:G/\phi^{-1}(U)\to G/U$ be the homomorphism induced by $\phi$, and note that $\widetilde\phi$ is injective.
%The proof is based on the following commuting diagram:
%\begin{equation*}
%\xymatrix{
%G \ar@/^1.5pc/[rr]^\chi \ar[d]^\phi \ar[r]^{\pi'} & \frac{G}{\phi^{-1}U} \ar[d]^{\widetilde\phi}\ar[r]^\eta & \mathbb T \\
%G \ar[r]^\pi \ar@/_3pc/[rru]_\theta & \frac{G}{U}\ar[ur]_{\xi} & 
%}
%\end{equation*}
%
%Let $\chi\in(\phi^{-1}(U))^\perp$. Then $\chi=\eta\circ\pi'$, where $\eta:G/\phi^{-1}(U)\to\mathbb T$. Since $\widetilde\phi$ is injective, and $\mathbb T$ is divisible, $\eta$ can be extended to $\xi:G/U\to \mathbb T$, i.e., $\eta=\xi\circ \widetilde\phi$. This gives $$\chi=\eta\circ \pi'=\xi\circ\widetilde\phi\circ\pi'=\xi\circ\pi\circ\phi,$$ and this shows that $\chi=\widehat\phi(\theta)$, where $\theta=\xi\circ\pi\in U^\perp$. This proves the inclusion $(\phi^{-1}(U))^\perp\subseteq\widehat\phi (U^\perp)$.
%
%Now let $\chi\in\widehat\phi (U^\perp)$. Then $\chi=\widehat\phi(\theta)=\theta\circ\phi$, where $\theta\in U^\perp$. So $\theta=\xi\circ\pi$, for some $\xi:G/U\to\mathbb T$. Take $\eta=\xi\circ\widetilde\phi$ (since $\widetilde\phi$ is injective we can think that $\eta=\xi\!\restriction_{G/\phi^{-1}(U)}$). Therefore, $$\chi=\theta\circ\phi=\xi\circ\pi\circ\phi=\xi\circ\widetilde\phi\circ\pi'=\eta\circ\pi'\in(\phi^{-1}(U))^\perp.$$ This prove the inclusion $(\phi^{-1}(U))^\perp\supseteq\widehat\phi (U^\perp)$.
%
%\smallskip
%The second assertion follows from the first one and from Lemma \ref{(v)}.
%\end{proof}
%
%\begin{lemma}\label{A/B}
%Let $G$ be a locally compact abelian group, and let $B\subseteq A$ be subgroups of $G$. Then $\widehat{A/B}$ is topologically isomorphic to $B^\perp/A^\perp$.
%\end{lemma}
%\begin{proof}
%Consider the embedding $A/B\to G/B$. By Pontryagin duality there exists an open continuous homomorphism $\widehat{G/B}\to \widehat{A/B}$. Since $\widehat{G/B}$ is topologically isomorphic to $B^\perp$, consider the composition $\varphi:B^\perp\to \widehat{A/B}$. Since $\ker\varphi=A^\perp$, the conclusion follows.
%\end{proof}
%
%\begin{proposition}\label{C_n-T_n}
%Let $G$ be a locally compact abelian group, let $\phi:G\to G$ be a continuous endomorphism, $U\in\mathcal B(G)$ and $K\in\mathcal B(\widehat G)$. Then, for every $n\in\N_+$,
%\begin{itemize}
%\item[(a)]$[U:C_n(\phi,U)]=[T_n(\widehat\phi,U^\perp):U^\perp].$
%\item[(b)]$[K^\top:C_n(\phi,K^\top)]=[T_n(\widehat\phi,K):K].$
%%$$[U:C_n(\phi,U)]=[T_n(\phi,U^\perp):U^\perp]$$
%\end{itemize}
%Hence,
%\begin{itemize}
%\item[(c)] $H_{top}(\phi,U)=H_{alg}(\widehat\phi,U^\perp);$
%\item[(d)] $H_{top}(\phi,K^\top)=H_{alg}(\widehat\phi,K).$
%\end{itemize}
%\end{proposition}
%\begin{proof}
%%Let $U\in\mathcal B(G)$. Then $K=U^\perp\in\mathcal B(\widehat G)$ by Lemma \ref{B(G)}. By Proposition \ref{perp}, $(\phi^{-n}(U))^\perp=(\widehat\phi)^n (K)$ for every $n\in\N$. Hence, $C_n(\phi,K)^\perp=T_n(\widehat \phi,K)$ for every $n\in\N_+$ by Lemma \ref{(v)}. It follows that
%%$[G:C_n(\phi,U)]=|\widehat{G/C_n(\phi,U)}|=|C_n(\phi,U)^\perp|=|T_n(\widehat\phi,K)|$ for every $n\in\N_+$.
%(a) Let $U\in\mathcal B(G)$. Then $U^\perp\in\mathcal B(\widehat G)$ by Lemma \ref{B(G)}. By Proposition \ref{perp}, $C_n(\phi,U)^\perp=T_n(\widehat \phi,U^\perp)$ for every $n\in\N_+$.
%Consider the following sequences:
%$$0\to C_n(\phi,U)\to U\to G$$
%and
%$$\widehat G=0^\perp\leftarrow T_n(\widehat\phi,U^\perp)=C_n(\phi,U)^\perp\leftarrow U^\perp\leftarrow 0=G^\perp.$$
%By Lemma \ref{A/B}, $\widehat{U/C_n(\phi,U)}\cong T_n(\widehat\phi,U^\perp)/U^\perp$. Since these quotients are finite, $U/C_n(\phi,U)\cong T_n(\widehat\phi,U^\perp)/U^\perp$, and this concludes the proof.
%
%\smallskip
%(b) follows from (a).
%
%\smallskip
%(c) and (d) follow respectively from (a) and (b). %, applying Proposition \ref{no-mu}.
%\end{proof}
%
%\begin{lemma}\label{cofinal}
%If $G$ is a totally disconnected locally compact abelian group such that $\widehat G$ is totally disconnected, then $\mathcal B(G)$ is cofinal in $\mathcal C(G)$.
%\end{lemma}
%\begin{proof}
%Let $K\in\mathcal C(G)$. Since $G$ is totally disconnected and locally compact abelian, $\mathcal B(G)$ is a base of the neighborhoods of $0$ in $G$ \cite{vD}. So there exists $U\in\mathcal B(G)$ such that $U\subseteq K$. Let $\pi:G\to G/U$ be the canonical projection. Since $G/U$ is discrete and torsion, and $\pi(K)$ is compact, it follows that $\pi(K)$ is finite. Then $F=\langle\pi(K)\rangle$ is a finite subgroup of $G/U$. Therefore, $\pi^{-1}(F)\in\mathcal B(G)$ and $\pi^{-1}(F)\supseteq K$. 
%\end{proof}
%
%\begin{theorem}
%Let $G$ be a totally disconnected locally compact abelian group such that $\widehat G$ is totally disconnected, and let $\phi:G\to G$ be a continuous endomorphism. Then $h_{top}(\phi)=h_{alg}(\widehat\phi)$. 
%\end{theorem}
%\begin{proof}
%Since $G$ is totally disconnected and locally compact abelian, $\mathcal B(G)$ is a base of the neighborhoods of $0$ in $G$ \cite{vD}.
%By Lemma \ref{base}(a), $h_{top}(\phi)=\sup\{H_{top}(\phi,U\}:U\in\mathcal B(G)\}$.
%
%Since $\widehat G$ is a totally disconnected locally compact abelian group such that $\widehat{\widehat G}$ is totally disconnected, being topologically isomorphic to $G$ by Pontryagin duality, it follows that $\mathcal B(\widehat G)$ is cofinal in $\mathcal C(\widehat G)$ in view of Lemma \ref{cofinal}. By Lemma \ref{base}(b), $h_{alg}(\widehat \phi)=\sup\{H_{alg}(\widehat\phi,K):K\in\mathcal B(\widehat G)\}$.
%
%By Proposition \ref{C_n-T_n}(c), $H_{top}(\phi,U)=H_{alg}(\widehat\phi,U^\perp)$ for every $U\in\mathcal B(G)$, so $h_{top}(\phi)\leq h_{alg}(\widehat \phi)$.
%By Proposition \ref{C_n-T_n}(d), $H_{top}(\phi,K^\top)=H_{alg}(\widehat\phi,K)$ for every $K\in\mathcal B(\widehat G)$, so $h_{top}(\phi)\geq h_{alg}(\widehat \phi)$, and this concludes the proof.
%\end{proof}



\begin{thebibliography}{00}

\small{
%\bibitem{A} D. Arnold, \emph{Finite Rank Torsion Free Abelian Groups and Rings }, L.N.M. n. 931, Springer, 1982.

\bibitem{AKM} R. L. Adler, A. G. Konheim, M. H. McAndrew, \emph{Topological entropy}, Trans. Amer. Math. Soc. 114 (1965) 309--319.

\bibitem{BDG} F. Berlai, D. Dikranjan, A. Giordano Bruno, \emph{Scale function vs Topological entropy}, Topology Appl. 160 (2013) 2314--2334.

\bibitem{B} R. Bowen, \emph{Entropy for group endomorphisms and homogeneous spaces}, Trans. Amer. Math. Soc. 153 (1971) 401--414.

%\bibitem{CH} Z. Chatzidakis, E. Hrushovski, \emph{An invariant for difference field extensions}, preprint \url{http://arxiv.org/abs/0902.0844v4}.

%\bibitem{DG} D. Dikranjan and A. Giordano Bruno, \emph{The Pinsker subgroup of an algebraic flow}, JPAA.

\bibitem{DG-BT2} D. Dikranjan, A. Giordano Bruno, \emph{The Bridge Theorem for totally disconnected LCA groups}, to appear in Topology Appl.

\bibitem{DG-BT} D. Dikranjan, A. Giordano Bruno, \emph{The connection between topological and algebraic entropy}, Topology Appl. 159 (13) (2012) 2980--2989.

\bibitem{DG-limitfree} D. Dikranjan, A. Giordano Bruno, \emph{Limit free computation of entropy}, Rend. Istit. Mat. Univ. Trieste 44 (2012) 297--312.

%\bibitem{DG2} D. Dikranjan, A. Giordano Bruno, \emph{Entropy on abelian groups}, preprint; arxiv:1007.0533.

\bibitem{DG-islam} D. Dikranjan, A. Giordano Bruno, \emph{Topological entropy and algebraic entropy for group endomorphisms}, Proceedings ICTA2011 Islamabad, Pakistan July 4-10 2011 Cambridge Scientific Publishers (2012) 133--214.

\bibitem{DGS} D. Dikranjan, A. Giordano Bruno, L. Salce, \emph{Adjoint algebraic entropy}, J. Algebra 324 (2010) 442--463.

\bibitem{DGSZ} D. Dikranjan, B. Goldsmith, L. Salce, P. Zanardo, \emph{Algebraic entropy of endomorphisms of abelian groups}, Trans. Amer. Math. Soc. {361} (2009) 3401--3434.

%\bibitem{DGZ} D. Dikranjan, K. Gong, P. Zanardo, \emph{Endomorphisms of abelian groups with small algebraic entropy}, submitted.

%\bibitem{DPS} D. Dikranjan, I. Prodanov and L. Stoyanov, \emph{Topological Groups: Characters,  Dualities  and  Minimal Group Topologies},  Pure and Applied Mathematics, Vol. \textbf{130}, Marcel Dekker Inc., New York-Basel, 1989.

%\bibitem{DSV} D. Dikranjan, M. Sanchis, S. Virili, {\em New and old facts about entropy on uniform spaces and topological groups}, to appear in Topology and Its Applications.

%\bibitem{Ward0}  G. Everest and T. Ward, \emph{Heights of polynomials and entropy in algebraic dynamics}, Universitext. Springer-Verlag London, Ltd., London, 1999. 

%\bibitem{fekete} M. Fekete, \emph{\"Uber die Verteilung der Wurzeln bei gewisser algebraichen Gleichungen mit ganzzahlingen Koeffizienten}, Math. Zeitschr. \textbf{17} (1923), 228--249.

%\bibitem{F} L. Fuchs, \emph{Infinite Abelian Groups}, Vol. I and II, Academic Press, 1970 and 1973.

%\bibitem{GBS} A. Giordano Bruno, L. Salce, \emph{A soft introduction to algebraic entropy}, to appear.

%\bibitem{GV} A. Giordano Bruno, S. Virili, \emph{The Algebraic Yuzvinski Formula}, preprint.

%\bibitem{Halmos} P. Halmos, \emph{On automorphisms of compact groups}, Bull. Amer. Math. Soc. {49} (1943) 619--624.

%\bibitem{Halmos2} P. Halmos, \emph{Finite-Dimensional Vector Spaces}, Undergraduate Texts in Mathematics, Springer-Verlag, New York - Heidelberg - Berlin, 1987.

\bibitem{HR}  E. Hewitt, K. A. Ross, \emph{Abstract harmonic analysis I}, Springer-Verlag, Berlin-Heidelberg-New York, 1963.

\bibitem{hood} B. M. Hood,  \emph{Topological entropy and uniform spaces}, J. London Math. Soc. 8 (2) (1974) 633--641. 

%\bibitem{KerLi} D. Kerr, H. Li,  \emph{Dynamical entropy in Banach spaces}, Invent. Math. \textbf{162} (2005), no. 3, 649--686. 

%\bibitem{Ki} B. Kitchens, \emph{Expansive dynamics of zero-dimensional groups}, Ergod. Th. and Dynam. Sys. 7 (1987) 249--261.

%\bibitem{Kr}  L. Kronecker, \emph{Zwei S\" atze \" uber Gleichungen mit ganzzahligen Coefficienten}, Jour. Reine Angew. Math. \textbf{53} (1857), 173--175.

%\bibitem{L} D. H. Lehmer, \emph{Factorization of certain cyclotomic functions}, Ann. of Math. \textbf{34} (1933), 461--469.

%\bibitem{LW} D. Lind and T. Ward, \emph{Automorphisms of solenoids and $p$-adic entropy}, Ergodic Theory Dynam. Systems \textbf{8} (1988) no. 3, 411-419.

\bibitem{P} L. S. Pontryagin, \emph{Topological Groups}, Gordon and Breach, New York, 1966.

%\bibitem{Pet} J. Peters, \emph{Entropy on discrete Abelian groups}, Adv. Math. {33} (1979) 1--13.

%\bibitem{Pet1}  J. Peters, \emph{Entropy of automorphisms on {L}.{C}.{A}. groups}, Pacific J. Math. 96 (2) (1981) 475--488.
         
\bibitem{St} L. N. Stoyanov, \emph{Uniqueness of topological entropy for endomorphisms on compact groups}, Boll. Un. Mat. Ital. B (7) 1 (3) (1987) 829--847.

\bibitem{vD} D. van Dantzig, \emph{Studien over topologische Algebra}, Dissertation, Amsterdam 1931.

%\bibitem{V} S. Virili, \emph{Entropy for endomorphisms of LCA groups}, to appear in Topology Appl.

%\bibitem{Wa} Walters, \emph{An Introduction to Ergodic Theory}, Springer-Verlag, New-York, 1982.

%\bibitem{Ward}  T. Ward, \emph{Group automorphisms with few and with many periodic points}, Proc. Amer. Math. Soc. \textbf{133} (2004), no.1, 91--96.

%\bibitem{WardLN}  T. Ward, \emph{Entropy of Compact Group Automorphisms}, on-line Lecture Notes (http://www.mth.uea.ac.uk/~h720/lecturenotes/).

\bibitem{W}  M. D. Weiss, \emph{Algebraic and other entropies of group endomorphisms}, Math. Systems Theory 8 (3) (1974/75) 243--248.

\bibitem{Willis} G. A. Willis, \emph{The structure of totally disconnected locally compact groups}, Math. Ann. 300 (2) (1994) 341--363. 

%\bibitem{Willis1} G. A. Willis, \emph{The number of prime factors of the scale function on a compactly generated group is finite}, Bull. London Math. Soc. \textbf{33} (2001) no. 2, 168--174.

\bibitem{Willis2} G. A. Willis, \emph{Further properties of the scale function on a totally disconnected group}, J. Algebra 237 (1) (2001) 142--164. 

%\bibitem{Willis3} G. Willis, {\em Tidy subgroups for commuting automorphisms of totally disconnected groups: An analogue of simultaneous triangularisation of matrices},  New York Journal of Mathematics,  10 (2004) 1--35.

%\bibitem{Willis4} G. A. Willis, \emph{The nub of an automorphism of a totally disconnected locally compact group}, submitted \url{http://arxiv.org/abs/1112.4239}.

%\bibitem{KepertWillis} A. Kepert, G. A. Willis, \emph{Scale functions and tree ends} J. Aust. Math. Soc. \textbf{70} (2001) no. 2, 273--292. 

%\bibitem{Y1} S. A. Yuzvinski, \emph{Calculation of the entropy of a group-endomorphism}, Sibirsk. Mat. \u Z. 8 (1967) 230--239.

\bibitem{Y} S. Yuzvinski, \emph{Metric properties of endomorphisms of compact groups}, Izv. Acad. Nauk SSSR, Ser. Mat. {29} (1965) 1295--1328 (in Russian). English Translation: Amer. Math. Soc. Transl. (2) {66} (1968) 63--98.

}

\end{thebibliography}

\end{document}